\section{Introduction}\label{sec:intro}
Several scientific theories of consciousness, and some general frameworks that have been brought to bear on it, describe the same open physical systems through different invariants. Integrated information theory (IIT) looks at how much a system's causal structure resists being cut~\cite{Oizumi2014}; global neuronal workspace theory (GNWT) looks at broadcast and global availability~\cite{DehaeneChangeux2011}; attention schema theory (AST) looks at an internal model of attention~\cite{GrazianoKastner2011}; the free-energy principle (FEP), an account of self-organizing systems, looks at Markov blankets and the statistics of the internal states that they shield~\cite{Friston2013}; and quantum Darwinism, an account of how classical records emerge, looks at redundant records in the environment~\cite{Zurek2009}. Category theory has already been used to study single theories and their relation to experience: the relation between consciousness and IIT~\cite{Northoff2019,Tsuchiya2016}; IIT itself, formulated in process theories, that is, in symmetric monoidal categories~\cite{TullKleiner2021}, on the basis of an axiomatic account of its mathematical structure~\cite{KleinerTull2021}; and graphical reasoning about conscious experience in process theories~\cite{SignorelliWangCoecke2021}. A further proposal is to unify several theories at once: to treat each one as a functor from a common category of open physical systems into its own category of descriptions, and to look for natural transformations, limits or bridge laws that connect them. This paper defines and audits a proposal of this kind, the candidate theory Q (Section~\ref{sec:obligacion}): a correspondence between open physical realizations and five components of phenomenal structure, $(c,q,u,\tau,\sigma)$. Q is the author's own proposal, so the audit is a self-audit, and it applies to Q the same standard that it applies to the proxies. Quantum-collapse proposals such as Orch-OR~\cite{HameroffPenrose2014}, which would enter Q only as additional postulates, are left aside: we define no proxy for them. The central open obligation of Q for the unity component $u$ is the unity sufficiency
\begin{equation}\label{eq:Su}
S_u:\qquad \Phi_Q(s)=\Phi_Q(t)\ \Longrightarrow\ u(s)=u(t),
\end{equation}
where $\Phi_Q$ is the physical descriptor of the theory and $s,t$ range over the original domain of realizations.

This paper is a formal and structural audit of that unification. Its aim is to delimit with mathematical precision which parts of it can be established, and which promises of unification fail, for explicit structural proxies of the theories, not for the theories themselves; the logical and combinatorial core of the argument is checked by computer in Lean~4, and its analytic core, three results on the Gaussian integration functional $J$ that serves as an IIT-type proxy, is checked in Lean~4 with the Mathlib library~\cite{Mathlib2020}. As a rule of method, the audit changes neither the descriptor, nor the domain, nor the components, nor the meaning of Q. We do not try to demonstrate Q, and we do not refute it, and the kind of evidence that supports each result is stated (Section~\ref{sec:status}). The paper pursues three aims, one per act.
\begin{itemize}
\item \emph{Act~I: negative delimitation} (Sections~\ref{sec:obligacion}--\ref{sec:debil}). First, the weak categorical thesis holds, but cheaply, and the naive strong unification fails: for the structural presence proxies of IIT type ($J>0$, with $J$ a Gaussian integration functional) and GNWT type (a broadcast hub), there is no translation between the two projections, their bridge square does not commute, and neither projection is monoidal into the additive reals. Second, the two proxies give opposite verdicts exactly on a band of connected systems without a hub, which is empty for $n\le3$ units and contains the path $P_n$ and the cycle $C_n$ for every $n\ge4$. Both vanish on systems with disconnected support, such as two disjoint edges, $2K_2$; since they vanish on tensor products of nonempty systems, they are strong monoidal into a zero-product monoid once the empty system is sent to an adjoined unit, whereas the extensive projections are additive. Third, we delimit the obligation $S_u$: we characterize exactly when a premise of the form ``the degree depends only on $\kappa$'' implies $S_u$, and we show that none of the candidate principles examined closes it without presupposing it, equating only isomorphic realizations, forcing the degree to be constant, changing the meaning of the degree, or changing the domain, as selecting the state $z=0$ of minimal surprisal does in a Gaussian family. With the degree a function of the coexperience structure, the banded segregation SEG that would be needed is non-degenerate $S_u$ itself, in testable form.
\item \emph{Act~II: validation in a Gaussian laboratory} (Section~\ref{sec:gauss}). In a two-unit linear-Gaussian family, exactly, and in four-unit networks, numerically, we confirm where the bridges commute (all systems with $n\le3$ units, and centralized systems such as the star $K_{1,3}$), where they break (the band, from $n\ge4$), and how sharp the min-cut bounds are.
\item \emph{Act~III: a positive route and a methodological guide} (Section~\ref{sec:via-futura}). Of the reformulations of $S_u$ examined, only intervention monotonicity escapes the collapse of Act~I. It is not $S_u$, nor an obligation of the original Q, but a declared weakening of its order clause B to a class of interventions on one bearer, and we have not shown that it is the only reformulation that escapes. We prove its exact frontier: on the intervention digraph, every degree that satisfies intervention monotonicity is constant on each strongly connected component, and degrees exist that vary across components. Taking $J$ as a stand-in descriptor, which is not $\Phi_Q$, we prove, and check in Lean with Mathlib, that, for at least two units, $J$ increases strictly under uniform scaling of a connected dynamics, and under scaling of the couplings with a fixed scalar self-coupling, and we show numerically that it can decrease if the fixed self-coupling is not scalar. From these results we draw a methodological guide for comparing theories by directed interventions within one system. The archive of the paper contains six Lean files without Mathlib and a Lean project with Mathlib of seven modules, all without \texttt{sorry}, and eleven symbolic and numerical scripts with their expected outputs.
\end{itemize}

\subsection{Main results, limitations and assumptions}\label{sec:summary}
\paragraph{Main results.} The status labels are explained in Section~\ref{sec:status}, and Table~\ref{tab:evidence} there lists the evidence for the principal results. Several of the formal results are elementary facts about kernels, factorizations and products: the entailment criterion (Theorem~\ref{thm:impl}), the translation criterion (\texttt{nat\_iff\_functional}), the pairing into the product (\texttt{pairing\_unique}) and the reconstruction criterion (Theorem~\ref{thm:reconstruccion}). We claim no novelty for them. What the paper adds is their application to the obligation $S_u$ of Q and to explicit proxies, the counterexamples, and the delimitation that follows.
\begin{enumerate}
\item[(R1)] \emph{Factorization premises.} Premises ``the degree depends only on $\kappa$'' imply $S_u$ if and only if $\ker\Phi_Q$ is contained in the equivalence relation they generate (Theorem~\ref{thm:impl}) \stL. If $\kappa$ is transversal to $\Phi_Q$ (Section~\ref{sec:intento}) and co-occurs connectedly with its levels, $S_u$ forces the degree to be constant (Theorem~\ref{thm:dicotomia}) \stL.
\item[(R2)] \emph{Candidate principles.} None of the principles examined in the available sources supplies an independent foundation excluding that two original realizations with the same $\Phi_Q$ differ in $u$ when $u$ is not constant (Sections~\ref{sec:principios}--\ref{sec:mejor}). The banded segregation SEG is not a new premise that would close $S_u$: with the degree a function of the coexperience structure, SEG with distinct values between groups is equivalent to non-degenerate $S_u$ (\texttt{segd\_iff}) \stL.
\item[(R3)] \emph{Categorical thesis.} The weak thesis holds as a scheme: every isomorphism-invariant assignment is a functor out of a symmetric monoidal groupoid of open systems (Proposition~\ref{prop:funtor}) \stA, and the projections pair uniquely into their product (\texttt{pairing\_unique}) \stL. It is cheap, because it holds for any family of invariants. The strong thesis fails for the structural presence bridges of IIT type ($J>0$, with $J$ a Gaussian integration functional) and GNWT type (a broadcast hub): their bridge square does not commute on the ring $C_4$ and on the path $P_4$, there is no translation from positive integration to broadcast, and joint projections have blind spots \stL; the reverse translation fails by the same criterion \stA. Lean without Mathlib encodes the IIT-type presence as connectivity of the support, and that this is $J>0$ is Lemma~\ref{lem:J}, \stA\ \stM\ for at least two units. The two presence projections vanish on tensor products of nonempty systems (Proposition~\ref{prop:monoidal} \stA; \texttt{tensor\_no\_universal} \stL), so they are not monoidal into the additive monoid $([0,+\infty],+,0)$ (Section~\ref{sec:proyecciones}). With the convention that the empty system goes to an adjoined unit and no nonempty system does, they are strong monoidal into a zero-product monoid (\texttt{strong\_monoidal\_zero}, in abstract form) \stL; without it they are not. The extensive projections (predictive information, and the mutual information of the FEP-type projection in the Gaussian case) are additive \stA.
\item[(R4)] \emph{The band.} The two presence bridges disagree exactly on connected systems without a broadcast hub: the Boolean form is compiled (\texttt{noncomm\_locus}, \texttt{commute\_off\_band}, \texttt{universal\_connected}) \stL, and the identification of $J>0$ with connectivity is Lemma~\ref{lem:J}, \stA\ \stM\ for at least two units. There is no such system with $n\le3$ units (\texttt{n2\_equiv}, \texttt{n3\_equiv}) \stL. The path $P_n$ is one for every $n\ge4$ (\texttt{path\_reach\_from\_zero}, \texttt{path\_no\_universal}) \stL, and so is the cycle $C_n$, whose maximum degree $2$ is smaller than $n-1$ \stA. For $n=4$ the band has $15$ labeled graphs (\texttt{n4\_count}) \stL, and for $n=5$ it has $472$ (exhaustive enumeration, \texttt{pistas\_numerico.py}) \stS.
\item[(R5)] \emph{Bounds.} For systems with at least two units, $J>0$ if and only if the undirected support is connected (Lemma~\ref{lem:J}, \texttt{QAudit.\allowbreak lemmaJ\_stationary}), and $J$ is sandwiched by the minimum cut $w_{\min}$, for each representative, since $w_{\min}$ is not an isomorphism invariant (Theorem~\ref{thm:cheeger}, \texttt{QAudit.\allowbreak sandwich}) \stA\ \stM. A broadcast hub of strength $c$ gives an explicit lower bound \stA. The lower bound of the sandwich is attained in a subfamily, so its constant cannot be improved \stS.
\item[(R6)] \emph{Gaussian laboratory.} In the two-unit families $F_0$ and $F_{\rm mem}$ of Section~\ref{sec:gauss} everything is exact ($156$ symbolic checks) \stS: $F_0$ and $F_0\otimes F_0$ realize the two levels on which the bridges commute; predictive information equals $J$ and the cut information vanishes; selecting the state of minimal surprisal, a free-energy-type selection, is trivial at the level of laws and discards almost every realization at the level of states and histories, since its minimizer, the mode, has probability zero; and on $F_0$ the premise IPS$_K$ (Definition~\ref{def:IPS}), that the degree depends only on the parameter $k=c^2s_2/s_1$ of the family, is equivalent to the sufficiency $S_J$ that it would be used to support. The four-unit networks are numerical \stN.
\item[(R7)] \emph{Positive route.} Exact covariance under perturbations and the order clause on all pairs collapse \stL. Intervention monotonicity does not: it forces the degree to be constant exactly on the strongly connected components of an intervention digraph (Theorem~\ref{thm:frontier}) \stL; in Lean the degrees take integer values, and the transcription to real values is not formalized (Section~\ref{sec:via-futura}). It is not $S_u$, nor an obligation of the original Q, but a declared weakening of the order clause B to a class of interventions on one bearer, and we have not shown that it is the only reformulation that escapes the collapse. In the band it can be instantiated through Proposition~\ref{prop:scaling} if $J$ is taken in place of $\Phi_Q$, which is not justified here: for systems with at least two units, $J(tA)$ increases strictly under uniform scaling for every connected support, and so does $J(aI+cW)$ in $c$ for a connected $W$ with a fixed scalar self-coupling $a$ (\texttt{QAuditScaling.\allowbreak prop\_scaling\_a}, \texttt{QAuditScaling.\allowbreak prop\_scaling\_b}) \stA\ \stM. Strengthening only the couplings while a non-scalar diagonal of self-couplings is held fixed can decrease $J$ \stN.
\end{enumerate}

\paragraph{Limitations.}
\begin{itemize}
\item \emph{Proxies.} $J$ is the Gaussian residual functional of geometric integrated information~\cite{Oizumi2016PNAS}, not the $\Phi$ of IIT~3.0~\cite{Oizumi2014}, and the broadcast hub is a structural stand-in for global availability, not the nonlinear ignition of GNWT. The disagreement on the band is a disagreement between these proxies computed on the same model and at the same grain, not between the theories as such.
\item \emph{Formalization.} Six Lean files use Lean~4 without Mathlib and check logical skeletons and finite combinatorics \stL. A Lean project with Mathlib checks three analytic statements about $J$ over real matrices \stM: Lemma~\ref{lem:J} and the min-cut sandwich (Theorem~\ref{thm:cheeger}) for $n\ge2$ units, the sandwich exactly as stated, with the extreme eigenvalues, and both parts of Proposition~\ref{prop:scaling}, the latter slightly more generally ((a) without assuming $\operatorname{spr}(A)<1$, (b) without assuming that $W$ has zero diagonal). In Lean the inner minimum over $\mathfrak B_\pi$ (Lemma~\ref{lem:J}) is an infimum, whose attainment is a separate checked theorem, and the stationary covariance is either given as data ($P\succ0$ with $P=APA^\top+E$, Section~\ref{sec:base}) or constructed as the series $\sum_{j\ge0}A^jE(A^j)^\top$, whose convergence is proved from $\operatorname{spr}(A)<1$, where $\operatorname{spr}$ is Mathlib's spectral radius of the complexified matrix. Not formalized: the conventions for fewer than two units (Remark~\ref{rem:small}), such as $J=+\infty$ for one unit; the derivative identity $Q'=AQ'A^\top+S$ for the matrix $Q(c)=c^2P(c)$ of the proof of Proposition~\ref{prop:scaling}(b) and identity~\eqref{eq:S-identity}, in whose place the Lean proof of Proposition~\ref{prop:scaling}(b) uses the inequality $Q(c_2)-Q(c_1)\succeq(c_2-c_1)^2E$, so that~\eqref{eq:S-identity} remains \stS\ for $3\times3$ matrices and \stA\ for general $n$; the uniqueness of the stationary covariance, which the Lean statements do not need; the vanishing of $J$ on tensor products (Proposition~\ref{prop:monoidal}); and Corollary~\ref{cor:band-scaling} as a single statement, which combines graph facts \stL\ with Proposition~\ref{prop:scaling}. The files without Mathlib and the Mathlib project are separate developments. The Lean statements of Section~\ref{sec:via-futura} on exact covariance, the order clause and the intervention digraph take degrees with integer values; their transcription to real values is not formalized. The hub bound (Theorem~\ref{thm:laxa}(d)), the additivity of the predictive information $I_{\rm pred}$ (Proposition~\ref{prop:monoidal}), Proposition~\ref{prop:funtor} and the other analytic statements remain paper proofs \stA, the Gaussian computations are symbolic checks in SymPy \stS, and the numerical results \stN\ are not certificates.
\item \emph{Self-audit.} Q is the author's own proposal. Every definition, hypothesis and result that the audit uses is stated in this paper (Section~\ref{sec:obligacion}), so the audit can be checked from the paper and its archive alone.
\item \emph{Q itself.} The original descriptor $\Phi_Q$, the original domain and the comparison relation for the degree are decisions of the author that have not been taken; the original obligations of Q remain open, and nothing here measures $u$.
\item \emph{Review.} All reviews of this work were carried out by AI agents (see the \hyperref[sec:ai]{declaration on the use of generative AI}); there has been no independent human review.
\end{itemize}

\paragraph{Assumptions.} The linear-Gaussian instance $\mathbf{Open}_G$ of the category of open systems (Section~\ref{sec:base}); the undirected-support reading of the broadcast hub; and the definitions of Section~\ref{sec:obligacion} as the framework audited.

\subsection{Status labels and verification}\label{sec:status}
The results of Section~\ref{sec:summary} and of Sections~\ref{sec:obligacion}--\ref{sec:via-futura} carry one or more of five labels; in Sections~\ref{sec:obligacion}--\ref{sec:debil} a label may be implicit, as explained below.
\begin{itemize}
\item \stL: checked by the kernel of Lean~4 (version~4.24.0, without Mathlib): logic and finite combinatorics. Each of the six Lean files compiles with exit code~0, contains no \texttt{sorry}, \texttt{admit}, \texttt{native\_decide}, \texttt{axiom} or \texttt{opaque} declarations, and depends at most on the standard axioms \texttt{propext}, \texttt{Classical.choice} and \texttt{Quot.sound}; the axioms of each theorem are listed in Appendix~\ref{sec:archivos}.
\item \stM: checked by the kernel of Lean~4 (version~4.24.0) with Mathlib~\cite{Mathlib2020} (v4.24.0) over real matrices. Each of the seven modules of the Lean project compiles with exit code~0 and contains no \texttt{sorry}, \texttt{admit}, \texttt{native\_decide}, \texttt{axiom} or \texttt{opaque} declarations, and every one of its $73$ \texttt{\#print axioms} reports is \texttt{[propext, Classical.choice, Quot.sound]} (Appendix~\ref{sec:archivos}). Where the Lean statement differs in scope from the statement of the paper, the difference is stated where the label is used and in Section~\ref{sec:summary} (\emph{Limitations}).
\item \stA: proved on paper in this paper. Alone, the label means that the result is not formalized; \stA\ \stM\ marks a result whose paper proof is kept and whose statement is also checked with Mathlib.
\item \stS: verified exactly by computer: symbolic identities and exact rational arithmetic in SymPy~1.14~\cite{SymPy2017} for the Gaussian model, or exact integer enumeration in Python.
\item \stN: numerical, without a certificate (local optimization with restarts, or high-precision root finding).
\end{itemize}
Where a result in Sections~\ref{sec:obligacion}--\ref{sec:debil} carries no explicit label, the label is implicit: a statement that cites a Lean name in \texttt{typewriter} font is checked by Lean in the abstract form stated next to it; Lean names whose namespace begins with \texttt{QAudit} belong to the Mathlib project and are \stM; the others belong to the files without Mathlib and are \stL. When a statement concerns $J$, the link between Lean's Boolean predicates and $J$ is Lemma~\ref{lem:J}, which is \stA\ \stM\ for at least two units; the files without Mathlib and the Mathlib project are separate developments, and this link is not itself a compiled statement. Lemma~\ref{lem:J} and Theorem~\ref{thm:cheeger} are \stA\ \stM\ for at least two units, Theorem~\ref{thm:laxa}(d) and the other analytic statements are \stA, and the values of $J$ in Table~\ref{tab:C-1} are \stN. Lean checks the logical form of a statement on abstract types and finite graphs and, with Mathlib, on real matrices; whether a physical or phenomenal premise holds is never decided by a compilation.

Table~\ref{tab:evidence} gives, for the principal results of Section~\ref{sec:summary}, the label and the Lean theorems or scripts that check each one; the other statements of Section~\ref{sec:summary} carry their labels and checks in the text. Appendix~\ref{sec:correspondencia} gives the complete correspondence, for every numbered result of the paper.

\begin{table}[htbp]\centering\footnotesize
\begin{tabular}{>{\raggedright\arraybackslash}p{0.06\textwidth}>{\raggedright\arraybackslash}p{0.42\textwidth}>{\raggedright\arraybackslash}p{0.08\textwidth}>{\raggedright\arraybackslash}p{0.32\textwidth}}
\toprule
Item & Main result & Status & Checked by\\
\midrule
(R1) & Entailment criterion (Theorem~\ref{thm:impl}) and co-occurrence dichotomy (Theorem~\ref{thm:dicotomia}) & \stL & \texttt{entails\_Su\_iff}, \texttt{connected\_forces\_constant}\\
(R2) & SEG with distinct values between groups is non-degenerate $S_u$ & \stL & \texttt{segd\_iff}\\
(R3) & Weak thesis: every isomorphism-invariant assignment is a functor (Proposition~\ref{prop:funtor}) & \stA & paper proof\\
(R3) & The bridge square fails on $C_4$ and $P_4$; no translation from positive integration to broadcast & \stL & \texttt{square\_fails}, \texttt{values}, \texttt{no\_translation}\\
(R3) & Strong monoidality into a zero-product monoid, in abstract form & \stL & \texttt{strong\_monoidal\_zero}\\
(R4) & Exact locus of disagreement of the two presence bridges (the band), in Boolean form; its link to $J>0$ is Lemma~\ref{lem:J} & \stL & \texttt{noncomm\_locus}, \texttt{commute\_off\_band}, \texttt{universal\_connected}\\
(R4) & The band is empty for $n\le3$, contains $P_n$ for every $n\ge4$, and has $15$ labeled graphs for $n=4$ & \stL & \texttt{n2\_equiv}, \texttt{n3\_equiv}, \texttt{path\_reach\_from\_zero}, \texttt{path\_no\_universal}, \texttt{n4\_count}\\
(R4) & The band has $472$ labeled graphs for $n=5$ & \stS & \path{pistas_numerico.py}\\
(R5) & For $n\ge2$ units, $J>0$ if and only if the support is connected (Lemma~\ref{lem:J}) & \stA\ \stM & \texttt{QAudit.\allowbreak lemmaJ\_stationary}\\
(R5) & For $n\ge2$ units, the min-cut sandwich (Theorem~\ref{thm:cheeger}) & \stA\ \stM & \texttt{QAudit.\allowbreak sandwich}\\
(R5) & The lower bound of the sandwich is attained in a subfamily & \stS & \path{gauss_conexiones.py}\\
(R6) & Gaussian laboratory: the $156$ exact checks on $F_0$, $F_{\rm mem}$ and $F_0\otimes F_0$ & \stS & \path{v12_familia_escala.py}, \path{v3_modelos_apuntados.py}, \path{v4_historias_ac.py}, \path{v5_rkp_cadenas.py}, \path{gauss_conexiones.py}\\
(R6) & The four-unit networks (Table~\ref{tab:networks}) & \stN & \path{pistas_numerico.py}, \path{J_conectividad.py}\\
(R7) & Collapse of exact covariance and of the order clause on all pairs & \stL & \texttt{exact\_covariance\_gives\_Su}, \texttt{order\_all\_pairs\_collapse}\\
(R7) & Exact frontier of intervention monotonicity (Theorem~\ref{thm:frontier}) & \stL & \texttt{im\_iff\_edges}, \texttt{constant\_on\_scc}, \texttt{nonconstant\_if\_not\_scc}\\
(R7) & For $n\ge2$ units and connected support, strict increase of $J(tA)$ in $t$ and of $J(aI+cW)$ in $c$ (Proposition~\ref{prop:scaling}) & \stA\ \stM & \texttt{QAuditScaling.\allowbreak prop\_scaling\_a}, \texttt{QAuditScaling.\allowbreak prop\_scaling\_b}\\
(R7) & $J$ can decrease when a non-scalar self-coupling is held fixed (Remark~\ref{rem:fixed-diagonal}) & \stN & \path{contraejemplo_diagonal_fija.py}\\
\bottomrule
\end{tabular}
\caption{Evidence for the main results of Section~\ref{sec:summary}: status label (Section~\ref{sec:status}) and the Lean theorems or scripts that check each one. Lean names written without their namespace are theorems of the six files without Mathlib; names written in the namespaces \texttt{QAudit} and \texttt{QAuditScaling} are theorems of the Lean project with Mathlib, with the scope stated in Section~\ref{sec:summary} (\emph{Limitations}); the Lean statements in the two \stL\ rows of (R7) take degrees with integer values. Appendix~\ref{sec:correspondencia} gives the complete correspondence.}\label{tab:evidence}
\end{table}
